\section{Previous Work}
\label{sec:previous_work}
% ============================================================
% MARKING CRITERION: "Literature Review" [40 marks]


\subsection{Dialogue, Vicarious Learning, and Educational Conversational Agents}
\label{sec:background}

For nearly two decades, spoken dialogue systems (SDS) have been
investigated in educational contexts, particularly for intelligent
tutoring \citep{litman2004itspoke, pon-barry2004scot} and language
learning \citep{seneff-etal-2007-spoken, bibauw2019dialogue}, and more
recently as conversational agents powered by large language models
\citep{Sydorenko2025}. Dialogue has been demonstrated to be a more effective modality than plain
text for knowledge retention \citep{ginns2013designing}, and merely
\emph{overhearing} dialogue between agents has itself been shown to
support learning better than listening to a monologue \citep{vicarious_learning}.

This phenomenon, coined vicarious learning, has been investigated
through "fictive dialogues": semi-scripted exchanges between virtual
agents generated from a source text, where the learner is a passive
listener rather than an active participant \citep{piwek2007t2d}.
\citet{fictive_dialogs} extend the paradigm to a classroom
language-learning Educational Conversational Agent (ECA) that
alternates between teacher and peer roles. They find that more
scripted agent and human-confederate exchanges improve comprehension, whereas
direct agent-to-learner questioning degrades it.
The multi-party configurations needed for designing such agentic
conversations remain a known open challenge for current systems
\citep{aylett2023you, skantze2021turntaking}, primarily because
realtime end-of-turn prediction, addressee detection and overlap
handling all degrade sharply once the dyadic speak/wait assumption
is dropped.

\subsection{LLM-based and Multi-Agent Tutors}
\label{sec:llm_tutors}

A recent systematic review of 36 educational chatbots finds that
teaching agents and peer agents together dominate deployed roles
\citep{kuhail2023interacting} --- a decomposition that maps directly
onto the teacher and question-asking student configuration the
present system adopts. Empirical evidence on the pedagogical effects
of the LLM-based variants of such systems, however, is only
beginning to accumulate. \citet{SpangenbergerPia2026CwaL}
evaluate chat-based conversation with an LLM in education for
sustainable development ($N=122$) and find that LLM chat engages both
affective and cognitive processes relevant to
learning, and yields knowledge gains comparable to reading the same
material as text. This establishes LLM chat as a viable instructional
medium rather than a pedagogical novelty. What unstructured dyadic
LLM chats still lack, however, is the curricular sequencing and
student modelling that drove the gains reported for earlier
conversational intelligent tutoring systems such as AutoTutor, which
was shown to produce learning gains on par with human tutors
\citep{graesser2016autotutor}.

A recent architectural response is to recover that pedagogical
structure by composing several specialised LLM agents rather than
relying on a single chatbot or hand-authored dialogue policy. The
WikiHowAgent system of \citet{pei-etal-2025conversational}
instantiates teacher, student, interaction-manager and evaluator
agents at scale; \citet{shi-etal-2025-educationq} apply a similar
triadic teacher--student--evaluator framework (EducationQ) to
benchmark teaching ability across 14 LLMs. Both systems place no
human in the interactive loop --- the human enters only as a
post-hoc judge. The present work adopts the same multi-agent
decomposition but narrows the student agent to a question-asking
role whose purpose is to elicit teacher exposition that a human
learner overhears or joins, over a realtime spoken channel rather
than text.

\subsection{Speech Interfaces for Educational Agents}
\label{sec:speech_interfaces}

The few LLM-era ECAs that do operate in the spoken modality
typically pair their dialogue manager with cloud-based speech I/O:
\citet{fictive_dialogs}, for instance, integrate Deepgram's hosted
ASR and TTS synthesis into a RASA dialogue agent.
\citet{Sydorenko2025} identify recognition error and end-to-end
latency as the two persistent bottlenecks for spoken educational
agents --- a concern echoed in the realtime turn-taking literature
\citep{skantze2021turntaking}. The present system adopts the same
two-component STT/TTS split but runs both locally (Whisper for STT,
Chatterbox for TTS), trading on-device compute for lower latency
and no external dependency in the realtime loop.

\subsection{Summary}
\label{sec:prev_work_summary}

Our system sits at the intersection of three traditions: the
dialogue-based intelligent tutoring tradition
\citep{litman2004itspoke, pon-barry2004scot} with its spoken,
pedagogically-structured interaction; the recent LLM-based tutor
wave \citep{SpangenbergerPia2026CwaL, pei-etal-2025conversational},
which contributes fluent generation and scalable multi-agent role
decomposition but operates in text-only, dyadic or
symbolically-routed modes; and the vicarious/fictive dialogue
paradigm
\citep{vicarious_learning, piwek2007t2d, fictive_dialogs}, which shows
that overhearing pedagogically structured dialogue between agents
improves learning. The system's central design move is to substitute
the scripted human confederate of \citet{fictive_dialogs} with a
second autonomous LLM agent, combining the teacher--learner agent
decomposition of \citet{pei-etal-2025conversational} with the
overhearing-learner stance of the fictive-dialogue tradition over a
realtime spoken channel. Closing this gap opens the door to emerging
application areas such as spoken classroom assistants,
accessibility-focused tutors for low-literacy learners, multi-party
language-learning environments, and healthcare communication training.
